\documentclass[11pt]{article}
\usepackage[margin=2cm]{geometry}
\usepackage[english,ngerman]{babel}
\usepackage{mathtools}
\usepackage{amsthm}
\usepackage[table]{xcolor}
\usepackage{siunitx}
\usepackage{graphicx}
\usepackage{booktabs}
\usepackage{tcolorbox}
\usepackage{listings}
\usepackage{tikz}
\usetikzlibrary{arrows.meta}
\usepackage{algorithm}
\usepackage{algpseudocode}
\usepackage[inline]{enumitem}
\usepackage{fancyhdr}
\usepackage[natbib=true]{biblatex}
\addbibresource{refs.bib}
\usepackage{glossaries}
\usepackage{acro}
\usepackage{marginnote}
\usepackage{hyperref}
\usepackage{cleveref}
\usepackage{qatools}

\newcommand{\vect}[1]{\mathbf{#1}}
\newcommand{\note}[2][red]{\textcolor{#1}{#2}}
\newtheorem{theorem}{Theorem}
\newcounter{step}
\definecolor{brandblue}{RGB}{20,60,160}
\newglossaryentry{latex}{name={LaTeX},description={a typesetting system}}
\DeclareAcronym{cpu}{short=CPU,long=central processing unit}

\begin{document}

\section{Introduction}\label{sec:intro}

Speed is $v^{2}$ and $x_{max}$.

\begin{equation}\label{eq:energy}
	E = mc^2
\end{equation}

\section{Results}\label{sec:results}

The vector $\vect{v}$ was measured.

\begin{theorem}\label{thm:main}
	A result.
\end{theorem}

\begin{figure}
	\includegraphics[width=0.4\textwidth]{figures/plot.png}
	\caption{A plot.}\label{fig:plot}
\end{figure}

Body text.


\printbibliography

\end{document}
